\item Durante periodos de gran actividad, el Sol tiene más manchas solares que de costumbre. Las manchas solares son más frías que el resto de la capa luminosa de la atmósfera del Sol (la fotosfera). Paradójicamente, la potencia de salida total del Sol activo no es menor que el promedio sino que es la misma o ligeramente mayor que el promedio. Elabore los detalles del siguiente modelo imperfecto de este fenómeno. Considere una mancha de la fotosfera con un área de $5.10 \times 10^{14}\text{ m}^2$. Su emisividad es 0.965. (a) Encuentre la potencia que radia si su temperatura es uniformemente de 5800 K, correspondiente al Sol quieto. (b) Para representar una mancha solar, suponga que 10.0\% del área del parche está a 4800 K y que el otro 90.0\% está a 5890 K. Encuentre la potencia de salida de la mancha. (c) Establezca cómo la respuesta al inciso (b) se compara con la respuesta al inciso (a). (d) Obtenga la temperatura promedio de la mancha. Note que esta temperatura más fría resulta en una alta presión de salida.
